\begin{pastexam}[label={px:ntu-106-dsa-03}]{臺大資工所 106 學年度 第 3 題}
（問答）

(15 points) A subsequence of a sequence $S$ is obtained by deleting zero or more elements from
$S$. Given two sequences, the longest common subsequence (LCS) problem is to find a subsequence
that is common to both sequences and its length is maximized. Complete the following pseudocode
for computing the length of an LCS and for delivering an LCS by filling \fbox{3A}, \fbox{3B} and
\fbox{3C}.

\smallskip
\noindent\rule{\linewidth}{0.4pt}
\begin{tabbing}
\hspace{1.5em}\=\hspace{1.5em}\=\hspace{1.5em}\=\hspace{1.5em}\=\kill
\textbf{Algorithm} \textsc{LCS\_length}($A=\langle a_1,a_2,\ldots,a_m\rangle$, $B=\langle b_1,b_2,\ldots,b_n\rangle$)\\
\textbf{begin}\\
\>\textbf{for} $i\gets 0$ \textbf{to} $m$ \textbf{do} $len[i,0]\gets 0$\\
\>\textbf{for} $j\gets 1$ \textbf{to} $n$ \textbf{do} $len[0,j]\gets 0$\\
\>\textbf{for} $i\gets 1$ \textbf{to} $m$ \textbf{do}\\
\>\>\textbf{for} $j\gets 1$ \textbf{to} $n$ \textbf{do}\\
\>\>\>\textbf{if} $a_i=b_j$ \textbf{then}\\
\>\>\>\>$len[i,j]\gets$ \fbox{3A}\\
\>\>\>\>$prev[i,j]\gets$``$\nwarrow$''\\
\>\>\>\textbf{else if} $len[i-1,j]\ge len[i,j-1]$ \textbf{then}\\
\>\>\>\>$len[i,j]\gets len[i-1,j]$\\
\>\>\>\>$prev[i,j]\gets$``$\uparrow$''\\
\>\>\>\textbf{else}\\
\>\>\>\>$len[i,j]\gets len[i,j-1]$\\
\>\>\>\>$prev[i,j]\gets$``$\leftarrow$''\\
\>\textbf{return} $len$ and $prev$\\
\textbf{end}\\[0.6em]
\textbf{Algorithm} \textsc{LCS\_output}($A=\langle a_1,a_2,\ldots,a_m\rangle$, $prev$, $i$, $j$)\\
\textbf{begin}\\
\>\textbf{if} $i=0$ \textbf{or} $j=0$ \textbf{then return}\\
\>\textbf{if} $prev[i,j]=$``$\nwarrow$'' \textbf{then}\\
\>\>\textsc{LCS\_output}(\fbox{3B})\\
\>\>\textbf{print} \fbox{3C}\\
\>\textbf{else if} $prev[i,j]=$``$\uparrow$'' \textbf{then} \textsc{LCS\_output}($A$, $prev$, $i-1$, $j$)\\
\>\textbf{else} \textsc{LCS\_output}($A$, $prev$, $i$, $j-1$)\\
\textbf{end}
\end{tabbing}
\vspace{-0.8em}
\noindent\rule{\linewidth}{0.4pt}
\end{pastexam}

\begin{pastsoln}
答案：
\begin{itemize}
\item \fbox{3A}：$len[i-1, j-1] + 1$
\item \fbox{3B}：$A,\ prev,\ i-1,\ j-1$
\item \fbox{3C}：$a_i$
\end{itemize}
最外層的呼叫是 \textsc{LCS\_output}($A$, $prev$, $m$, $n$)。

這兩個演算法就是資結書虛擬碼~18.6（\proc{LCS-Length}）和虛擬碼~18.7（\proc{Print-LCS}），只是把
箭頭表叫做 $prev$。

\textbf{3A}。 $len[i, j]$ 是 $a_1 \ldots a_i$ 和 $b_1 \ldots b_j$ 的 LCS 長度。 $a_i = b_j$ 時，任取 $a_1 \ldots a_i$ 和 $b_1 \ldots b_j$ 的一個 LCS $Z = z_1 \ldots z_k$。 如果
$z_k \ne a_i$，$Z$ 後面再接上 $a_i$ 仍然是共同子序列，比 $Z$ 長，矛盾；所以 $z_k = a_i = b_j$，而且
可以讓 $z_k$ 配到 $a_i$、$b_j$ 這兩個位置。 去掉 $z_k$，前面的 $z_1 \ldots z_{k-1}$ 是 $a_1 \ldots a_{i-1}$
和 $b_1 \ldots b_{j-1}$ 的共同子序列，而且必須是最長的，否則換成更長的，再接上 $a_i$，會得到比 $Z$
更長的共同子序列（資結書定理~18.2 的第一種情形）。 所以
\[
len[i, j] = len[i-1, j-1] + 1 ,
\]
箭頭記成「$\nwarrow$」，表示這一格是從左上角那一格加一個字元來的。 另外兩種情形（$a_i \ne b_j$）
已經寫在題目裡：LCS 至少丟掉 $a_i$ 或 $b_j$ 其中一個，取上面一格和左邊一格的較大者。

\textbf{3B、3C}。 \textsc{LCS\_output}($A$, $prev$, $i$, $j$) 印出 $a_1 \ldots a_i$ 和 $b_1 \ldots b_j$ 的
一個 LCS。 箭頭是「$\nwarrow$」時，這一格的 LCS 是「$(i-1, j-1)$ 那一格的 LCS，後面接上
$a_i$」。 要依原本的順序印出來，先遞迴印前面的部分，再印最後這個字元：
\[
\textsc{LCS\_output}(A, prev, i-1, j-1); \quad \textbf{print } a_i .
\]
印 $a_i$ 和印 $b_j$ 一樣（兩者相等），但函式只拿到 $A$，所以只能寫 $a_i$。 先印再遞迴的話，順序會
反過來。

\textbf{小例子}。 $A = \mathtt{ABCBDAB}$、$B = \mathtt{BDCABA}$。 依上面
填好的演算法，$len[7, 6] = 4$，從 $(7, 6)$ 沿著箭頭往回走，碰到「$\nwarrow$」的格子依序是
$(6,6)\,\mathtt{A}$、$(4,5)\,\mathtt{B}$、$(3,3)\,\mathtt{C}$、$(2,1)\,\mathtt{B}$（由後往前），
遞迴回來時由前往後印出 \texttt{BCBA}。

\textbf{複雜度}。 \textsc{LCS\_length} 填 $(m+1)(n+1)$ 格，每格 $O(1)$，共 $\Theta(mn)$ 時間和空間。
\textsc{LCS\_output} 每次呼叫 $i$ 或 $j$ 至少減一，最多 $m + n$ 次呼叫，$O(m + n)$。
\end{pastsoln}
